\section{Семинар 4. Отношения, отображения и образы множеств}
\noeqref{eq:rel_q_def, eq:map_uniq_def, eq:func_idemp, eq:im_preim_def}
\noeqref{eq:im_cap, eq:im_cup, eq:im_diff, eq:preim_cap, eq:shift_map}

\begin{remark}[Обзор темы]
    На четвёртом семинаре закрепляются навыки проверки аксиом частичного 
    порядка и отношения эквивалентности, исследуется связь между 
    теоретико-множественным и функциональным определениями отображения, 
    доказываются фундаментальные свойства прямых образов и полных 
    прообразов множеств относительно операций объединения, пересечения и 
    разности, а также изучаются свойства композиции и инъективности.
\end{remark}

\subsection{Задачи на отношения порядка и эквивалентности}

\begin{problem}
    Доказать, что множество всех подмножеств (булеан) $\powerset(M)$ 
    произвольного множества $M$ частично упорядочено отношением нестрогого 
    включения $\subseteq$.
\end{problem}

\begin{proof}[Решение]~\\
    \begin{quote}
        Проверим выполнение трёх аксиом нестрогого частичного порядка для 
        бинарного отношения $\subseteq$ на множестве $\powerset(M)$:
        \begin{enumerate}
            \item \textbf{Рефлексивность ($\forall A \in \powerset(M) : 
            A \subseteq A$):} 
            По определению включения условие $A \subseteq A$ означает, что 
            $\forall x \, (x \in A \implies x \in A)$. Импликация вида 
            $P \implies P$ является тавтологией, поэтому любое множество 
            является подмножеством самого себя.
            \item \textbf{Антисимметричность ($\forall A, B \in \powerset(M) : 
            (A \subseteq B \land B \subseteq A) \implies A = B$):} 
            Если каждый элемент множества $A$ принадлежит $B$ и каждый 
            элемент множества $B$ принадлежит $A$, то множества $A$ и $B$ 
            состоят из одних и тех же элементов, что по аксиоме объёмности 
            (определению равенства множеств) означает $A = B$.
            \item \textbf{Транзитивность ($\forall A, B, C \in \powerset(M) : 
            (A \subseteq B \land B \subseteq C) \implies A \subseteq C$):} 
            Пусть $x \in A$ --- произвольный элемент. Так как $A \subseteq B$, 
            то $x \in B$. Так как $B \subseteq C$, из $x \in B$ следует 
            $x \in C$. Следовательно, $\forall x \, (x \in A \implies x \in C)$, 
            то есть $A \subseteq C$.
        \end{enumerate}
        Все три свойства выполнены, следовательно, пара 
        $\bigl(\powerset(M), \subseteq\bigr)$ образует частично 
        упорядоченное множество (ЧУМ).
    \end{quote}
\end{proof}

\begin{problem}
    На множестве вещественных чисел $\R$ задано бинарное отношение $Q$ по 
    правилу:
    \begin{equation} \label{eq:rel_q_def}
        \alpha Q \beta \iff (\alpha - \beta) \in \Q \tag{1}
    \end{equation}
    Доказать, что $Q$ является отношением эквивалентности на $\R$, и 
    описать его классы эквивалентности.
\end{problem}

\begin{proof}[Решение]~\\
    \begin{quote}
        Проверим рефлексивность, симметричность и транзитивность отношения 
        $Q$, опираясь на замкнутость поля рациональных чисел $\Q$ 
        относительно арифметических операций:
        \begin{enumerate}
            \item \textbf{Рефлексивность:} Для любого $\alpha \in \R$ имеем 
            $\alpha - \alpha = 0 = \frac{0}{1} \in \Q$, следовательно, 
            $\alpha Q \alpha$.
            \item \textbf{Симметричность:} Пусть $\alpha Q \beta$, то есть 
            $\alpha - \beta = r \in \Q$. Тогда противоположное число 
            $\beta - \alpha = -(\alpha - \beta) = -r$ также является 
            рациональным ($-r \in \Q$), откуда $\beta Q \alpha$.
            \item \textbf{Транзитивность:} Пусть $\alpha Q \beta$ и 
            $\beta Q \gamma$. Это означает, что $\alpha - \beta = \omega_1 \in \Q$ 
            и $\beta - \gamma = \omega_2 \in \Q$. Представим разность 
            $\alpha - \gamma$ в виде телескопической суммы:
            \[
                \alpha - \gamma = (\alpha - \beta) + (\beta - \gamma) = 
                \omega_1 + \omega_2.
            \]
            Поскольку сумма двух рациональных чисел рациональна 
            ($\omega_1 + \omega_2 \in \Q$), получаем $\alpha Q \gamma$.
        \end{enumerate}
        Таким образом, $Q$ --- отношение эквивалентности. Класс 
        эквивалентности произвольного числа $\alpha \in \R$ имеет вид 
        сдвинутого множества рациональных чисел:
        \[
            \cl_Q(\alpha) = \{\beta \in \R \mid \alpha - \beta \in \Q\} = 
            \{\alpha + r \mid r \in \Q\} = \alpha + \Q.
        \]
        Фактор-множество $\R / Q = \R / \Q$ состоит из несчётного 
        (континуального) семейства счётных классов эквивалентности (именно 
        выбор по одному представителю из каждого такого класса на отрезке 
        $[0, 1]$ позволяет построить классическое неизмеримое по Лебегу 
        множество Витали).
    \end{quote}
\end{proof}

\subsection{Отображения как бинарные отношения}

\begin{definition}
    Бинарное отношение $f \subseteq X \times Y$ называется \textdef{полным 
    отображением} (или \textdef{всюду определённой функцией}) из $X$ в $Y$, 
    если каждому элементу $x \in X$ соответствует ровно один элемент $y \in Y$:
    \begin{equation} \label{eq:map_uniq_def}
        \forall x \in X \ \exists! y \in Y : x f y \tag{2}
    \end{equation}
    При этом вместо инфиксной записи $x f y$ (или $(x, y) \in f$) 
    используется функциональная запись $f(x) = y$. Если же для каждого 
    $x \in X$ существует \emph{не более одного} $y \in Y$, такое отношение 
    называется \textdef{частичным (функциональным) отображением}.
\end{definition}

\begin{problem}
    Пусть $f \subseteq A \times A$ --- отображение на множестве $A$. 
    Исследовать, при каком условии отношение $f$ является отношением 
    нестрогого частичного порядка на $A$, и выяснить связь этого условия 
    с тождеством идемпотентности:
    \begin{equation} \label{eq:func_idemp}
        f\bigl(f(x)\bigr) = f(x) \tag{3}
    \end{equation}
\end{problem}

\begin{proof}[Решение и разбор случаев]~\\
    \begin{quote}
        На записи $x f y \iff f(x) = y$ проверим, к чему приводят аксиомы 
        частичного порядка (рефлексивность, антисимметричность, 
        транзитивность). Разберём два случая: когда $f$ понимается как 
        произвольное функциональное отношение (частичное отображение) и 
        когда $f$ всюду определено на $A$.

        \textbf{1. Анализ транзитивности и антисимметричности (связь с 
        идемпотентностью \eqref{eq:func_idemp}).}\\
        Пусть $x \in \dom f$ и $f(x) = y$ (то есть $x f y$).
        \begin{itemize}
            \item($\implies$) Предположим, что $f$ --- частичный порядок 
            (или хотя бы транзитивное и рефлексивное на своём образе 
            функциональное отношение). Если $y \in \dom f$ и $f(y) = z$ 
            (то есть $y f z$), то по свойству транзитивности из $x f y$ и 
            $y f z$ следует $x f z$, что означает $f(x) = z$. Но в силу 
            однозначности отображения из $f(x) = y$ и $f(x) = z$ получаем 
            $y = z$ (при $y \neq z$ возникает противоречие с функциональностью). 
            Следовательно:
            \[
                f\bigl(f(x)\bigr) = f(y) = z = y = f(x),
            \]
            то есть выполняется тождество идемпотентности 
            \eqref{eq:func_idemp}.
            \item ($\impliedby$) Обратно, пусть для функционального 
            отношения $f$ на всей его области определения выполняется 
            $f(f(x)) = f(x)$. Проверим антисимметричность и транзитивность:
            \begin{enumerate}
                \item \emph{Антисимметричность:} Пусть $x f y$ и $y f x$, 
                то есть $f(x) = y$ и $f(y) = x$. Тогда:
                \[
                    x = f(y) = f\bigl(f(x)\bigr) = f(x) = y \implies x = y.
                \]
                \item \emph{Транзитивность:} Пусть $x f y$ и $y f z$, то 
                есть $f(x) = y$ и $f(y) = z$. Тогда:
                \[
                    z = f(y) = f\bigl(f(x)\bigr) = f(x) \implies x f z.
                \]
            \end{enumerate}
        \end{itemize}

        \textbf{2. Учёт рефлексивности на всём множестве $A$.}\\
        Аксиома рефлексивности требует, чтобы для \emph{каждого} $x \in A$ 
        выполнялось соотношение $x f x$, то есть $f(x) = x$. Так как для 
        отображения образ каждого элемента $x$ единственен, условие 
        $f(x) = x$ при всех $x \in A$ однозначно задаёт тождественное 
        отображение $f = I_A = e_A$ (для которого, очевидно, 
        $f(f(x)) = x = f(x)$). Таким образом:
        \begin{itemize}
            \item Функциональное отношение $f \subseteq A^2$ транзитивно и 
            антисимметрично (а также рефлексивно на множестве своих 
            значений $\im f$) тогда и только тогда, когда оно идемпотентно: 
            $f(f(x)) = f(x)$.
            \item Полное отображение $f: A \to A$ является нестрогим 
            частичным порядком на всём множестве $A$ тогда и только тогда, 
            когда оно тождественно: $\forall x \in A : f(x) = x$.
        \end{itemize}
    \end{quote}
\end{proof}

\subsection{Образ и полный прообраз множества при отображении}

\begin{definition}
    Пусть задано отображение $f: X \to Y$, а также подмножества 
    $A \subseteq X$ и $B \subseteq Y$.
    \begin{itemize}
        \item \textdef{Образом (прямым образом)} множества $A$ при 
        отображении $f$ называется множество всех значений функции $f$ на 
        элементах из $A$:
        \item \textdef{Полным прообразом} множества $B$ при отображении $f$ 
        называется множество всех элементов из $X$, образы которых попадают 
        в $B$:
    \end{itemize}
    \begin{equation} \label{eq:im_preim_def}
        f(A) = \{f(x) \mid x \in A\} = \{y \in Y \mid \exists x \in A : f(x) = y\},
        \qquad
        f^{-1}(B) = \{x \in X \mid f(x) \in B\}. \tag{4}
    \end{equation}
\end{definition}

\begin{remark}
    Обозначение полного прообраза $f^{-1}(B)$ в формуле \eqref{eq:im_preim_def} 
    имеет смысл для \emph{любого} отображения $f: X \to Y$ и не требует 
    существования обратного отображения $f^{-1}: Y \to X$.
\end{remark}

\begin{problem}
    Пусть дано отображение $f: X \to Y$. Доказать следующие соотношения для 
    подмножеств области определения и области значений:
    \begin{align}
        \text{a) } & f(A \cap B) \subseteq f(A) \cap f(B) 
        && \forall A, B \subseteq X; \label{eq:im_cap} \tag{5} \\
        \text{b) } & f(A \cup B) = f(A) \cup f(B) 
        && \forall A, B \subseteq X; \label{eq:im_cup} \tag{6} \\
        \text{c) } & f(A) \setminus f(B) \subseteq f(A \setminus B) 
        && \forall A, B \subseteq X; \label{eq:im_diff} \tag{7} \\
        \text{d) } & f^{-1}(A \cap B) = f^{-1}(A) \cap f^{-1}(B) 
        && \forall A, B \subseteq Y. \label{eq:preim_cap} \tag{8}
    \end{align}
\end{problem}

\begin{proof}[Решение]~\\
    \begin{quote}
        Докажем последовательно каждое из четырёх утверждений.

        \textbf{Пункт a): $f(A \cap B) \subseteq f(A) \cap f(B)$.}\\
        Пусть $y \in f(A \cap B)$. По определению образа 
        \eqref{eq:im_preim_def} существует элемент $x \in A \cap B$, такой 
        что $f(x) = y$. Так как $x \in A \cap B$, то одновременно $x \in A$ 
        и $x \in B$.
        \begin{itemize}
            \item Из $x \in A$ и $y = f(x)$ следует $y \in f(A)$;
            \item Из $x \in B$ и $y = f(x)$ следует $y \in f(B)$.
        \end{itemize}
        Значит, $y \in f(A) \cap f(B)$, что доказывает включение 
        \eqref{eq:im_cap}.

        \textbf{Пункт b): $f(A \cup B) = f(A) \cup f(B)$.}\\
        Докажем взаимное включение левой и правой частей:
        \begin{enumerate}
            \item Включение «$\subseteq$»: пусть $y \in f(A \cup B)$. Тогда 
            найдётся $x \in A \cup B$, для которого $f(x) = y$. По 
            определению объединения $x \in A$ или $x \in B$. В первом 
            случае $y = f(x) \in f(A)$, во втором --- $y = f(x) \in f(B)$. 
            Следовательно, $y \in f(A) \cup f(B)$.
            \item Включение «$\supseteq$»: пусть $y \in f(A) \cup f(B)$, 
            тогда $y \in f(A)$ или $y \in f(B)$. Если $y \in f(A)$, то 
            существует $x \in A \subseteq A \cup B$, такой что $f(x) = y$, 
            откуда $y \in f(A \cup B)$. Аналогично, если $y \in f(B)$, то 
            существует $x \in B \subseteq A \cup B$ с условием $f(x) = y$, 
            и снова $y \in f(A \cup B)$.
        \end{enumerate}
        Из двух встречных включений следует равенство \eqref{eq:im_cup}.

        \textbf{Пункт c): $f(A) \setminus f(B) \subseteq f(A \setminus B)$.}\\
        Пусть $y \in f(A) \setminus f(B)$. Это означает, что $y \in f(A)$ и 
        $y \notin f(B)$. Так как $y \in f(A)$, существует элемент $x \in A$, 
        для которого $f(x) = y$. Покажем, что этот элемент $x$ не может 
        принадлежать $B$: действительно, если бы выполнялось $x \in B$, то 
        его образ $y = f(x)$ принадлежал бы $f(B)$, что противоречит 
        условию $y \notin f(B)$. Следовательно, $x \in A$ и $x \notin B$, 
        то есть $x \in A \setminus B$, откуда $y = f(x) \in f(A \setminus B)$.

        \textbf{Пункт d): $f^{-1}(A \cap B) = f^{-1}(A) \cap f^{-1}(B)$.}\\
        Докажем равенство \eqref{eq:preim_cap} в обе стороны:
        \begin{enumerate}
            \item Включение «$\subseteq$»: пусть $x \in f^{-1}(A \cap B)$. 
            По определению полного прообраза \eqref{eq:im_preim_def} имеем 
            $f(x) \in A \cap B$, то есть $f(x) \in A$ и $f(x) \in B$. Но 
            $f(x) \in A \implies x \in f^{-1}(A)$, а 
            $f(x) \in B \implies x \in f^{-1}(B)$. Следовательно, 
            $x \in f^{-1}(A) \cap f^{-1}(B)$.
            \item Включение «$\supseteq$»: пусть 
            $x \in f^{-1}(A) \cap f^{-1}(B)$, тогда $x \in f^{-1}(A)$ и 
            $x \in f^{-1}(B)$. По определению прообраза это означает, что 
            $f(x) \in A$ и $f(x) \in B$, то есть $f(x) \in A \cap B$. 
            Отсюда по определению получаем $x \in f^{-1}(A \cap B)$.
        \end{enumerate}
    \end{quote}
\end{proof}

\begin{remark}[\textbf{Почему в пунктах a) и c) нет равенства}]
    Если отображение $f$ не является инъективным, то включения 
    \eqref{eq:im_cap} и \eqref{eq:im_diff} могут быть строгими. Приведём 
    контрпример: пусть $f: \R \to \R$, $f(x) = x^2$, и выбраны множества 
    $A = \{-1\}$, $B = \{1\}$. Тогда:
    \begin{itemize}
        \item $A \cap B = \varnothing \implies f(A \cap B) = \varnothing$, 
        однако $f(A) \cap f(B) = \{1\} \cap \{1\} = \{1\} \neq \varnothing$.
        \item $A \setminus B = \{-1\} \implies f(A \setminus B) = \{1\}$, 
        однако $f(A) \setminus f(B) = \{1\} \setminus \{1\} = \varnothing \neq \{1\}$.
    \end{itemize}
    Равенства $f(A \cap B) = f(A) \cap f(B)$ и 
    $f(A \setminus B) = f(A) \setminus f(B)$ выполняются для любых 
    подмножеств $A, B \subseteq X$ тогда и только тогда, когда отображение 
    $f$ \textbf{инъективно}. В то же время полный прообраз $f^{-1}$ 
    сохраняет все теоретико-множественные операции (объединение, 
    пересечение, разность, дополнение) для \emph{любого} отображения $f$.
\end{remark}

\subsection{Композиция и свойства конкретных отображений}

\begin{problem}
    Пусть бинарные отношения $g \subseteq X \times Y$ и 
    $f \subseteq Y \times Z$ являются отображениями ($g: X \to Y$ и 
    $f: Y \to Z$). Доказать строго по определению \eqref{eq:map_uniq_def}, 
    что их композиция $f \circ g \subseteq X \times Z$ также является 
    отображением.
\end{problem}

\begin{proof}[Решение]~\\
    \begin{quote}
        Напомним, что по определению композиции отношений пара 
        $(x, z) \in f \circ g$ тогда и только тогда, когда существует 
        промежуточный элемент $y \in Y$, такой что $(x, y) \in g$ и 
        $(y, z) \in f$. Проверим для $f \circ g$ существование и 
        единственность образа каждого элемента $x \in X$:
        \begin{enumerate}
            \item \textbf{Всюду определённость (существование образа):} 
            Возьмём произвольный элемент $x \in X$. Так как $g: X \to Y$ --- 
            отображение, существует элемент $y \in Y$, такой что $x g y$. 
            Для этого элемента $y \in Y$, поскольку $f: Y \to Z$ --- 
            отображение, существует элемент $z \in Z$, такой что $y f z$. 
            Следовательно, для любого $x \in X$ существует $z \in Z$, для 
            которого $x (f \circ g) z$.
            \item \textbf{Функциональность (единственность образа):} 
            Пусть для некоторого $x \in X$ выполняются соотношения 
            $x (f \circ g) z_1$ и $x (f \circ g) z_2$. По определению 
            композиции существуют такие элементы $y_1, y_2 \in Y$, что:
            \[
                (x g y_1 \ \text{и} \ y_1 f z_1) \qquad \text{и} \qquad
                (x g y_2 \ \text{и} \ y_2 f z_2).
            \]
            Так как $g$ --- отображение, из условий $x g y_1$ и $x g y_2$ 
            в силу единственности образа следует $y_1 = y_2 = y$. Подставляя 
            $y_1 = y_2 = y$ во вторые соотношения, получаем $y f z_1$ и 
            $y f z_2$. Так как $f$ --- отображение, отсюда следует 
            $z_1 = z_2$.
        \end{enumerate}
        Таким образом, $\forall x \in X \ \exists! z \in Z : x (f \circ g) z$, 
        то есть $f \circ g: X \to Z$ является отображением.
    \end{quote}
\end{proof}

\begin{problem}
    Исследовать на инъективность, сюръективность и биективность отображение 
    сдвига $f: \N \to \N$ (где $\N = \{1, 2, 3, \dots\}$), заданное формулой:
    \begin{equation} \label{eq:shift_map}
        f(n) = n + 1 \tag{9}
    \end{equation}
\end{problem}

\begin{proof}[Решение]~\\
    \begin{quote}
        Проверим свойства инъективности и сюръективности по определению:
        \begin{enumerate}
            \item \textbf{Инъективность:} Пусть $x, y \in \N$ и $x \neq y$. 
            Прибавляя к обеим частям неравенства единицу, получаем 
            $x + 1 \neq y + 1$, то есть $f(x) \neq f(y)$ (или, 
            эквивалентно, $f(x) = f(y) \implies x + 1 = y + 1 \implies x = y$). 
            Следовательно, отображение $f$ \textbf{инъективно}.
            \item \textbf{Сюръективность:} Найдём образ множества $\N$ при 
            отображении $f$:
            \[
                \im f = \{n + 1 \mid n \in \N\} = \{2, 3, 4, \dots\} = 
                \N \setminus \{1\}.
            \]
            У наименьшего натурального числа $y = 1 \in \N$ не существует 
            прообраза в $\N$, так как уравнение $n + 1 = 1$ даёт корень 
            $n = 0 \notin \N$ (если же считать $0 \in \N$, то прообраза в 
            $\N$ не будет у элемента $y = 0$, поскольку $n = -1 \notin \N$). 
            Следовательно, отображение $f$ \textbf{не сюръективно} 
            (а значит, и \textbf{не биективно}).
        \end{enumerate}
    \end{quote}
\end{proof}

\begin{remark}
    Рассмотренное отображение \eqref{eq:shift_map} демонстрирует 
    характерное свойство бесконечных множеств (определение бесконечности по 
    Дедекинду): бесконечное множество $\N$ может быть инъективно (и 
    биективно на образ) отображено на своё собственное подмножество 
    $\N \setminus \{1\} \subset \N$. Для конечных множеств это невозможно: 
    любое инъективное отображение конечного множества в себя обязательно 
    является сюръективным и биективным.
\end{remark}